\documentclass[a4paper]{article}
\usepackage{tikz}
\pgfmathsetmacro{\l}{5}
\title{Making a Tangram}
\author{Disha Kuzhively}
\date{}
\begin{document}
\maketitle
\begin{enumerate}
    \item[Step 1.] Take a sheet of paper and cut out a square.
    \item[Step 2.] Cut two congruent triangles from this square.
    \item[Step 3.] Cut one of the triangles again into two congruent triangles.
    \item[Step 4.] Of the remaining big triangle, find a mid-point of the longest side and mark it with a press.
    \item[Step 5.] Take the vertex opposite to the longest side and superimpose it on the mid-point of the longest side and cut the small triangle formed.
    \item[Step 6.] Fold the trapezium to get two congruent trapeziums
    \item[Step 7.]
\end{enumerate}
\begin{figure}
    \centering
    \begin{tikzpicture}
        \draw (0,0) rectangle (\l,\l);
    \end{tikzpicture}
    \caption{Step 1}
\end{figure}
\begin{figure}
    \centering
    \begin{tikzpicture}
        \draw (0,0) rectangle (\l,\l);
        \draw[dashed] (0,0) -- (\l,\l);
    \end{tikzpicture}
    \begin{tikzpicture}
        \draw (0,0) -- (\l,\l) -- (0,\l) -- cycle;
        \draw[xshift=10] (0,0) -- (\l,\l) -- (\l,0) -- cycle;
    \end{tikzpicture}
    \caption{Step 2}
\end{figure}
\begin{figure}
    \centering
    \begin{tikzpicture}
        \draw (0,0) -- (\l,\l) -- (0,\l) -- cycle;
        \draw[dashed] (0,\l) -- (\l/2,\l/2);
        \draw[xshift=10] (0,0) -- (\l,\l) -- (\l,0) -- cycle;
    \end{tikzpicture}
    \begin{tikzpicture}
        \draw (0,0) -- (\l/2,\l/2) -- (0,\l) -- cycle;
        \draw[yshift=10] (\l/2,\l/2) -- (0,\l) -- (\l,\l) -- cycle;
        \draw[xshift=10] (0,0) -- (\l,\l) -- (\l,0) -- cycle;
    \end{tikzpicture}
    \caption{Step 3}
\end{figure}
\begin{figure}
    \centering
    \begin{tikzpicture}
        \draw (0,0) -- (\l/2,\l/2) -- (0,\l) -- cycle;
        \draw[yshift=10] (\l/2,\l/2) -- (0,\l) -- (\l,\l) -- cycle;
        \draw[xshift=10] (0,0) -- (\l,\l) -- (\l,0) -- cycle;
        \draw[xshift=10,dashed] (\l/2,\l/2) --++ (-45:1) -- cycle;
        \filldraw[gray,xshift=10] (0,0) circle [radius=2pt]
                         (\l,\l) circle [radius=2pt];   
    \end{tikzpicture}
    \caption{Step 4}
\end{figure}
\begin{figure}
    \centering
    \begin{tikzpicture}
        \draw (0,0) -- (\l/2,\l/2) -- (0,\l) -- cycle;
        \draw[yshift=10] (\l/2,\l/2) -- (0,\l) -- (\l,\l) -- cycle;
        \draw[xshift=10] (0,0) -- (\l,\l) -- (\l,\l/2) -- (\l/2,\l/2) -- (\l/2,0) -- cycle;
        \filldraw[gray,xshift=10] (\l/2,0) --  (\l/2,\l/2) -- (\l,\l/2) -- cycle;   
    \end{tikzpicture}
        \begin{tikzpicture}
        \draw (0,0) -- (\l/2,\l/2) -- (0,\l) -- cycle;
        \draw[yshift=10] (\l/2,\l/2) -- (0,\l) -- (\l,\l) -- cycle;
        \draw[xshift=10] (0,0) -- (\l,\l) -- (\l,\l/2) -- (\l/2,0) -- cycle;
        \draw[xshift=20] (\l/2,0) -- (\l,\l/2) -- (\l,0) -- cycle;
    \end{tikzpicture}
    \caption{Step 5}
\end{figure}
\begin{figure}
    \centering
        \begin{tikzpicture}
        \draw (0,0) -- (\l/2,\l/2) -- (0,\l) -- cycle;
        \draw[yshift=10] (\l/2,\l/2) -- (0,\l) -- (\l,\l) -- cycle;
        \draw[xshift=10] (0,0) -- (\l,\l) -- (\l,\l/2) -- (\l/2,0) -- cycle;
        \draw[xshift=20] (\l/2,0) -- (\l,\l/2) -- (\l,0) -- cycle;
        \draw[xshift=20] (\l/2,\l/2) -- ()
    \end{tikzpicture}
    \caption{Step 6}
\end{figure}
\end{document}